
\documentclass{article}
\usepackage{colortbl}
\usepackage{makecell}
\usepackage{multirow}
\usepackage{supertabular}

\begin{document}

\newcounter{utterance}

\centering \large Interaction Transcript for game `matchit\_ascii', experiment `similar\_grid\_1', episode 2 with Qwen3.6{-}35B{-}A3B{-}FP8{-}without{-}reasoning.
\vspace{24pt}

{ \footnotesize  \setcounter{utterance}{1}
\setlength{\tabcolsep}{0pt}
\begin{supertabular}{c@{$\;$}|p{.15\linewidth}@{}p{.15\linewidth}p{.15\linewidth}p{.15\linewidth}p{.15\linewidth}p{.15\linewidth}}
    \# & $\;$A & \multicolumn{4}{c}{Game Master} & $\;\:$B\\
    \hline

    \theutterance \stepcounter{utterance}  
    & & \multicolumn{4}{p{0.6\linewidth}}{
        \cellcolor[rgb]{0.9,0.9,0.9}{
            \makecell[{{p{\linewidth}}}]{
                \texttt{\tiny{[P1$\langle$GM]}}
                \texttt{Zúčastňujete sa spoločnej hry na hádanie.} \\
\texttt{Cieľom je zistiť, či je táto mriežka a iný obrázok, ktorý vidím len ja, rovnaké.} \\
\texttt{Najskôr opíšte svoju mriežku. Potom poskytnem svoj opis a môžeme si navzájom klásť otázky o mriežkach, aby sme zistili, či sú rovnaké. Teraz začnite svoj krátky opis mriežky s "POPIS:", za ktorým nasleduje opis. Nepridávajte nič iné. Toto je vaša mriežka:} \\
\texttt{X X X X X} \\
\texttt{X ▢ ▢ ▢ ▢} \\
\texttt{X X X ▢ ▢} \\
\texttt{X ▢ ▢ ▢ ▢} \\
\texttt{X ▢ ▢ ▢ ▢} \\
            }
        }
    }
    & \\ \\

    \theutterance \stepcounter{utterance}  
    & \multicolumn{4}{p{0.6\linewidth}}{
        \cellcolor[rgb]{0.9,0.9,0.9}{
            \makecell[{{p{\linewidth}}}]{
                \texttt{\tiny{[P1$\rangle$GM]}}
                \texttt{POPIS: Mriežka 5x5 pozostávajúca z X a prázdnych štvorcov (▢). Prvý stĺpec je úplne vyplnený znakom X. Druhý riadok má na začiatku X a nasleduje štyri prázdne štvorce. Tretí riadok má na začiatku tri X a potom dva prázdne štvorce. Štvrtý a piaty riadok majú na začiatku X a nasledujú štyri prázdne štvorce.} \\
            }
        }
    }
    & & \\ \\

    \theutterance \stepcounter{utterance}  
    & & & \multicolumn{2}{p{0.3\linewidth}}{
        \cellcolor[rgb]{0.9,0.9,0.9}{
            \makecell[{{p{\linewidth}}}]{
                \texttt{\tiny{[GM$|$GM]}}
                \texttt{continue} \\
            }
        }
    }
    & & \\ \\

    \theutterance \stepcounter{utterance}  
    & & \multicolumn{4}{p{0.6\linewidth}}{
        \cellcolor[rgb]{0.9,0.9,0.9}{
            \makecell[{{p{\linewidth}}}]{
                \texttt{\tiny{[GM$\rangle$P2]}}
                \texttt{Zúčastňujete sa spoločnej hry na hádanie.} \\
\texttt{Cieľom je zistiť, či je táto mriežka a iný obrázok, ktorý vidím len ja, rovnaké.} \\
\texttt{Najskôr opíšte svoju mriežku. Potom poskytnem svoj opis a môžeme si navzájom klásť otázky o mriežkach, aby sme zistili, či sú rovnaké. Teraz začnite svoj krátky opis mriežky s "POPIS:", za ktorým nasleduje opis. Nepridávajte nič iné. Toto je vaša mriežka:} \\
\texttt{X X X X X} \\
\texttt{▢ ▢ ▢ ▢ X} \\
\texttt{▢ ▢ X X X} \\
\texttt{▢ ▢ ▢ ▢ X} \\
\texttt{▢ ▢ ▢ ▢ X} \\
            }
        }
    }
    & \\ \\

    \theutterance \stepcounter{utterance}  
    & & & \multicolumn{4}{p{0.6\linewidth}}{
        \cellcolor[rgb]{0.9,0.9,0.9}{
            \makecell[{{p{\linewidth}}}]{
                \texttt{\tiny{[GM$\langle$P2]}}
                \texttt{POPIS: Mriežka má 5 riadkov a 5 stĺpcov. Prvý riadok obsahuje päť krížikov (X). Druhý riadok má štyri štvorce (▢) a na konci jeden krížik (X). Tretí riadok má dva štvorce (▢) a tri krížiky (X). Štvrtý a piaty riadok sú identické a obsahujú štyri štvorce (▢) a na konci jeden krížik (X).} \\
            }
        }
    }
     \\ \\

    \theutterance \stepcounter{utterance}  
    & & & \multicolumn{2}{p{0.3\linewidth}}{
        \cellcolor[rgb]{0.9,0.9,0.9}{
            \makecell[{{p{\linewidth}}}]{
                \texttt{\tiny{[GM$|$GM]}}
                \texttt{continue} \\
            }
        }
    }
    & & \\ \\

    \theutterance \stepcounter{utterance}  
    & & \multicolumn{4}{p{0.6\linewidth}}{
        \cellcolor[rgb]{0.9,0.9,0.9}{
            \makecell[{{p{\linewidth}}}]{
                \texttt{\tiny{[GM$\rangle$P2]}}
                \texttt{Toto je môj POPISPOPIS: Mriežka 5x5 pozostávajúca z X a prázdnych štvorcov (▢). Prvý stĺpec je úplne vyplnený znakom X. Druhý riadok má na začiatku X a nasleduje štyri prázdne štvorce. Tretí riadok má na začiatku tri X a potom dva prázdne štvorce. Štvrtý a piaty riadok majú na začiatku X a nasledujú štyri prázdne štvorce.} \\
\texttt{Teraz položte otázku, aby ste zistili nové aspekty mojej mriežky, ktoré sa môžu líšiť od vašej mriežky. Začnite s "OTÁZKA:" a nepridávajte nič iné.} \\
            }
        }
    }
    & \\ \\

    \theutterance \stepcounter{utterance}  
    & & & \multicolumn{4}{p{0.6\linewidth}}{
        \cellcolor[rgb]{0.9,0.9,0.9}{
            \makecell[{{p{\linewidth}}}]{
                \texttt{\tiny{[GM$\langle$P2]}}
                \texttt{OTÁZKA: V treťom riadku je podľa vášho popisu na začiatku tri X a potom dva prázdne štvorce. Môžete potvrdiť, či je poradie v tomto riadku presne X, X, X, ▢, ▢?} \\
            }
        }
    }
     \\ \\

    \theutterance \stepcounter{utterance}  
    & & & \multicolumn{2}{p{0.3\linewidth}}{
        \cellcolor[rgb]{0.9,0.9,0.9}{
            \makecell[{{p{\linewidth}}}]{
                \texttt{\tiny{[GM$|$GM]}}
                \texttt{continue} \\
            }
        }
    }
    & & \\ \\

    \theutterance \stepcounter{utterance}  
    & & \multicolumn{4}{p{0.6\linewidth}}{
        \cellcolor[rgb]{0.9,0.9,0.9}{
            \makecell[{{p{\linewidth}}}]{
                \texttt{\tiny{[P1$\langle$GM]}}
                \texttt{Toto je môj POPISPOPIS: Mriežka má 5 riadkov a 5 stĺpcov. Prvý riadok obsahuje päť krížikov (X). Druhý riadok má štyri štvorce (▢) a na konci jeden krížik (X). Tretí riadok má dva štvorce (▢) a tri krížiky (X). Štvrtý a piaty riadok sú identické a obsahujú štyri štvorce (▢) a na konci jeden krížik (X).} \\
\texttt{OTÁZKA: V treťom riadku je podľa vášho popisu na začiatku tri X a potom dva prázdne štvorce. Môžete potvrdiť, či je poradie v tomto riadku presne X, X, X, ▢, ▢?Svoju odpoveď začnite s "ODPOVEĎ:" a nepridávajte nič iné.} \\
            }
        }
    }
    & \\ \\

    \theutterance \stepcounter{utterance}  
    & \multicolumn{4}{p{0.6\linewidth}}{
        \cellcolor[rgb]{0.9,0.9,0.9}{
            \makecell[{{p{\linewidth}}}]{
                \texttt{\tiny{[P1$\rangle$GM]}}
                \texttt{ODPOVEĎ: Áno, v treťom riadku je poradie presne X, X, X, ▢, ▢.} \\
            }
        }
    }
    & & \\ \\

    \theutterance \stepcounter{utterance}  
    & & & \multicolumn{2}{p{0.3\linewidth}}{
        \cellcolor[rgb]{0.9,0.9,0.9}{
            \makecell[{{p{\linewidth}}}]{
                \texttt{\tiny{[GM$|$GM]}}
                \texttt{continue} \\
            }
        }
    }
    & & \\ \\

    \theutterance \stepcounter{utterance}  
    & & \multicolumn{4}{p{0.6\linewidth}}{
        \cellcolor[rgb]{0.9,0.9,0.9}{
            \makecell[{{p{\linewidth}}}]{
                \texttt{\tiny{[P1$\langle$GM]}}
                \texttt{Teraz položte otázku, aby ste zistili nové aspekty mojej mriežky, ktoré sa môžu líšiť od vašej mriežky. Začnite s "OTÁZKA:" a nepridávajte nič iné.} \\
            }
        }
    }
    & \\ \\

    \theutterance \stepcounter{utterance}  
    & \multicolumn{4}{p{0.6\linewidth}}{
        \cellcolor[rgb]{0.9,0.9,0.9}{
            \makecell[{{p{\linewidth}}}]{
                \texttt{\tiny{[P1$\rangle$GM]}}
                \texttt{OTÁZKA: V druhom riadku je podľa vášho popisu na konci jeden krížik (X). Môžete potvrdiť, či je poradie v tomto riadku presne ▢, ▢, ▢, ▢, X?} \\
            }
        }
    }
    & & \\ \\

    \theutterance \stepcounter{utterance}  
    & & & \multicolumn{2}{p{0.3\linewidth}}{
        \cellcolor[rgb]{0.9,0.9,0.9}{
            \makecell[{{p{\linewidth}}}]{
                \texttt{\tiny{[GM$|$GM]}}
                \texttt{continue} \\
            }
        }
    }
    & & \\ \\

    \theutterance \stepcounter{utterance}  
    & & \multicolumn{4}{p{0.6\linewidth}}{
        \cellcolor[rgb]{0.9,0.9,0.9}{
            \makecell[{{p{\linewidth}}}]{
                \texttt{\tiny{[GM$\rangle$P2]}}
                \texttt{ODPOVEĎ: Áno, v treťom riadku je poradie presne X, X, X, ▢, ▢.} \\
\texttt{OTÁZKA: V druhom riadku je podľa vášho popisu na konci jeden krížik (X). Môžete potvrdiť, či je poradie v tomto riadku presne ▢, ▢, ▢, ▢, X?Svoju odpoveď začnite s "ODPOVEĎ:" a nepridávajte nič iné.} \\
            }
        }
    }
    & \\ \\

    \theutterance \stepcounter{utterance}  
    & & & \multicolumn{4}{p{0.6\linewidth}}{
        \cellcolor[rgb]{0.9,0.9,0.9}{
            \makecell[{{p{\linewidth}}}]{
                \texttt{\tiny{[GM$\langle$P2]}}
                \texttt{ODPOVEĎ: Áno, v druhom riadku je poradie presne ▢, ▢, ▢, ▢, X.} \\
            }
        }
    }
     \\ \\

    \theutterance \stepcounter{utterance}  
    & & & \multicolumn{2}{p{0.3\linewidth}}{
        \cellcolor[rgb]{0.9,0.9,0.9}{
            \makecell[{{p{\linewidth}}}]{
                \texttt{\tiny{[GM$|$GM]}}
                \texttt{continue} \\
            }
        }
    }
    & & \\ \\

    \theutterance \stepcounter{utterance}  
    & & \multicolumn{4}{p{0.6\linewidth}}{
        \cellcolor[rgb]{0.9,0.9,0.9}{
            \makecell[{{p{\linewidth}}}]{
                \texttt{\tiny{[GM$\rangle$P2]}}
                \texttt{Teraz položte otázku, aby ste zistili nové aspekty mojej mriežky, ktoré sa môžu líšiť od vašej mriežky. Začnite s "OTÁZKA:" a nepridávajte nič iné.} \\
            }
        }
    }
    & \\ \\

    \theutterance \stepcounter{utterance}  
    & & & \multicolumn{4}{p{0.6\linewidth}}{
        \cellcolor[rgb]{0.9,0.9,0.9}{
            \makecell[{{p{\linewidth}}}]{
                \texttt{\tiny{[GM$\langle$P2]}}
                \texttt{OTÁZKA: V štvrtom a piatom riadku ste uviedli, že majú na začiatku X a nasledujú štyri prázdne štvorce. Môžete potvrdiť, či je poradie v týchto riadkoch presne X, ▢, ▢, ▢, ▢?} \\
            }
        }
    }
     \\ \\

    \theutterance \stepcounter{utterance}  
    & & & \multicolumn{2}{p{0.3\linewidth}}{
        \cellcolor[rgb]{0.9,0.9,0.9}{
            \makecell[{{p{\linewidth}}}]{
                \texttt{\tiny{[GM$|$GM]}}
                \texttt{continue} \\
            }
        }
    }
    & & \\ \\

    \theutterance \stepcounter{utterance}  
    & & \multicolumn{4}{p{0.6\linewidth}}{
        \cellcolor[rgb]{0.9,0.9,0.9}{
            \makecell[{{p{\linewidth}}}]{
                \texttt{\tiny{[P1$\langle$GM]}}
                \texttt{ODPOVEĎ: Áno, v druhom riadku je poradie presne ▢, ▢, ▢, ▢, X.} \\
\texttt{OTÁZKA: V štvrtom a piatom riadku ste uviedli, že majú na začiatku X a nasledujú štyri prázdne štvorce. Môžete potvrdiť, či je poradie v týchto riadkoch presne X, ▢, ▢, ▢, ▢?Svoju odpoveď začnite s "ODPOVEĎ:" a nepridávajte nič iné.} \\
            }
        }
    }
    & \\ \\

    \theutterance \stepcounter{utterance}  
    & \multicolumn{4}{p{0.6\linewidth}}{
        \cellcolor[rgb]{0.9,0.9,0.9}{
            \makecell[{{p{\linewidth}}}]{
                \texttt{\tiny{[P1$\rangle$GM]}}
                \texttt{ODPOVEĎ: Áno, v štvrtom a piatom riadku je poradie presne X, ▢, ▢, ▢, ▢.} \\
            }
        }
    }
    & & \\ \\

    \theutterance \stepcounter{utterance}  
    & & & \multicolumn{2}{p{0.3\linewidth}}{
        \cellcolor[rgb]{0.9,0.9,0.9}{
            \makecell[{{p{\linewidth}}}]{
                \texttt{\tiny{[GM$|$GM]}}
                \texttt{continue} \\
            }
        }
    }
    & & \\ \\

    \theutterance \stepcounter{utterance}  
    & & \multicolumn{4}{p{0.6\linewidth}}{
        \cellcolor[rgb]{0.9,0.9,0.9}{
            \makecell[{{p{\linewidth}}}]{
                \texttt{\tiny{[P1$\langle$GM]}}
                \texttt{Teraz položte otázku, aby ste zistili nové aspekty mojej mriežky, ktoré sa môžu líšiť od vašej mriežky. Začnite s "OTÁZKA:" a nepridávajte nič iné.} \\
            }
        }
    }
    & \\ \\

    \theutterance \stepcounter{utterance}  
    & \multicolumn{4}{p{0.6\linewidth}}{
        \cellcolor[rgb]{0.9,0.9,0.9}{
            \makecell[{{p{\linewidth}}}]{
                \texttt{\tiny{[P1$\rangle$GM]}}
                \texttt{OTÁZKA: V prvom riadku ste uviedli, že obsahuje päť krížikov (X). Môžete potvrdiť, či je poradie v tomto riadku presne X, X, X, X, X?} \\
            }
        }
    }
    & & \\ \\

    \theutterance \stepcounter{utterance}  
    & & & \multicolumn{2}{p{0.3\linewidth}}{
        \cellcolor[rgb]{0.9,0.9,0.9}{
            \makecell[{{p{\linewidth}}}]{
                \texttt{\tiny{[GM$|$GM]}}
                \texttt{continue} \\
            }
        }
    }
    & & \\ \\

    \theutterance \stepcounter{utterance}  
    & & \multicolumn{4}{p{0.6\linewidth}}{
        \cellcolor[rgb]{0.9,0.9,0.9}{
            \makecell[{{p{\linewidth}}}]{
                \texttt{\tiny{[GM$\rangle$P2]}}
                \texttt{ODPOVEĎ: Áno, v štvrtom a piatom riadku je poradie presne X, ▢, ▢, ▢, ▢.} \\
\texttt{OTÁZKA: V prvom riadku ste uviedli, že obsahuje päť krížikov (X). Môžete potvrdiť, či je poradie v tomto riadku presne X, X, X, X, X?Svoju odpoveď začnite s "ODPOVEĎ:" a nepridávajte nič iné.} \\
            }
        }
    }
    & \\ \\

    \theutterance \stepcounter{utterance}  
    & & & \multicolumn{4}{p{0.6\linewidth}}{
        \cellcolor[rgb]{0.9,0.9,0.9}{
            \makecell[{{p{\linewidth}}}]{
                \texttt{\tiny{[GM$\langle$P2]}}
                \texttt{ODPOVEĎ: Áno, v prvom riadku je poradie presne X, X, X, X, X.} \\
            }
        }
    }
     \\ \\

    \theutterance \stepcounter{utterance}  
    & & & \multicolumn{2}{p{0.3\linewidth}}{
        \cellcolor[rgb]{0.9,0.9,0.9}{
            \makecell[{{p{\linewidth}}}]{
                \texttt{\tiny{[GM$|$GM]}}
                \texttt{continue} \\
            }
        }
    }
    & & \\ \\

    \theutterance \stepcounter{utterance}  
    & & \multicolumn{4}{p{0.6\linewidth}}{
        \cellcolor[rgb]{0.9,0.9,0.9}{
            \makecell[{{p{\linewidth}}}]{
                \texttt{\tiny{[GM$\rangle$P2]}}
                \texttt{Teraz položte otázku, aby ste zistili nové aspekty mojej mriežky, ktoré sa môžu líšiť od vašej mriežky. Začnite s "OTÁZKA:" a nepridávajte nič iné.} \\
            }
        }
    }
    & \\ \\

    \theutterance \stepcounter{utterance}  
    & & & \multicolumn{4}{p{0.6\linewidth}}{
        \cellcolor[rgb]{0.9,0.9,0.9}{
            \makecell[{{p{\linewidth}}}]{
                \texttt{\tiny{[GM$\langle$P2]}}
                \texttt{OTÁZKA: V druhom stĺpci je podľa vášho popisu v prvom riadku X, v druhom riadku prázdny štvorec, v treťom riadku X, v štvrtom a piatom riadku prázdny štvorec. Môžete potvrdiť, či je poradie v druhom stĺpci presne X, ▢, X, ▢, ▢?} \\
            }
        }
    }
     \\ \\

    \theutterance \stepcounter{utterance}  
    & & & \multicolumn{2}{p{0.3\linewidth}}{
        \cellcolor[rgb]{0.9,0.9,0.9}{
            \makecell[{{p{\linewidth}}}]{
                \texttt{\tiny{[GM$|$GM]}}
                \texttt{continue} \\
            }
        }
    }
    & & \\ \\

    \theutterance \stepcounter{utterance}  
    & & \multicolumn{4}{p{0.6\linewidth}}{
        \cellcolor[rgb]{0.9,0.9,0.9}{
            \makecell[{{p{\linewidth}}}]{
                \texttt{\tiny{[P1$\langle$GM]}}
                \texttt{ODPOVEĎ: Áno, v prvom riadku je poradie presne X, X, X, X, X.} \\
\texttt{OTÁZKA: V druhom stĺpci je podľa vášho popisu v prvom riadku X, v druhom riadku prázdny štvorec, v treťom riadku X, v štvrtom a piatom riadku prázdny štvorec. Môžete potvrdiť, či je poradie v druhom stĺpci presne X, ▢, X, ▢, ▢?Svoju odpoveď začnite s "ODPOVEĎ:" a nepridávajte nič iné.} \\
            }
        }
    }
    & \\ \\

    \theutterance \stepcounter{utterance}  
    & \multicolumn{4}{p{0.6\linewidth}}{
        \cellcolor[rgb]{0.9,0.9,0.9}{
            \makecell[{{p{\linewidth}}}]{
                \texttt{\tiny{[P1$\rangle$GM]}}
                \texttt{ODPOVEĎ: Áno, v druhom stĺpci je poradie presne X, ▢, X, ▢, ▢.} \\
            }
        }
    }
    & & \\ \\

    \theutterance \stepcounter{utterance}  
    & & & \multicolumn{2}{p{0.3\linewidth}}{
        \cellcolor[rgb]{0.9,0.9,0.9}{
            \makecell[{{p{\linewidth}}}]{
                \texttt{\tiny{[GM$|$GM]}}
                \texttt{continue} \\
            }
        }
    }
    & & \\ \\

    \theutterance \stepcounter{utterance}  
    & & \multicolumn{4}{p{0.6\linewidth}}{
        \cellcolor[rgb]{0.9,0.9,0.9}{
            \makecell[{{p{\linewidth}}}]{
                \texttt{\tiny{[P1$\langle$GM]}}
                \texttt{Teraz položte otázku, aby ste zistili nové aspekty mojej mriežky, ktoré sa môžu líšiť od vašej mriežky. Začnite s "OTÁZKA:" a nepridávajte nič iné.} \\
            }
        }
    }
    & \\ \\

    \theutterance \stepcounter{utterance}  
    & \multicolumn{4}{p{0.6\linewidth}}{
        \cellcolor[rgb]{0.9,0.9,0.9}{
            \makecell[{{p{\linewidth}}}]{
                \texttt{\tiny{[P1$\rangle$GM]}}
                \texttt{OTÁZKA: V treťom stĺpci je podľa vášho popisu v prvom riadku X, v druhom riadku prázdny štvorec, v treťom riadku X, v štvrtom a piatom riadku prázdny štvorec. Môžete potvrdiť, či je poradie v treťom stĺpci presne X, ▢, X, ▢, ▢?} \\
            }
        }
    }
    & & \\ \\

    \theutterance \stepcounter{utterance}  
    & & & \multicolumn{2}{p{0.3\linewidth}}{
        \cellcolor[rgb]{0.9,0.9,0.9}{
            \makecell[{{p{\linewidth}}}]{
                \texttt{\tiny{[GM$|$GM]}}
                \texttt{continue} \\
            }
        }
    }
    & & \\ \\

    \theutterance \stepcounter{utterance}  
    & & \multicolumn{4}{p{0.6\linewidth}}{
        \cellcolor[rgb]{0.9,0.9,0.9}{
            \makecell[{{p{\linewidth}}}]{
                \texttt{\tiny{[GM$\rangle$P2]}}
                \texttt{ODPOVEĎ: Áno, v druhom stĺpci je poradie presne X, ▢, X, ▢, ▢.} \\
\texttt{OTÁZKA: V treťom stĺpci je podľa vášho popisu v prvom riadku X, v druhom riadku prázdny štvorec, v treťom riadku X, v štvrtom a piatom riadku prázdny štvorec. Môžete potvrdiť, či je poradie v treťom stĺpci presne X, ▢, X, ▢, ▢?Svoju odpoveď začnite s "ODPOVEĎ:" a nepridávajte nič iné.} \\
            }
        }
    }
    & \\ \\

    \theutterance \stepcounter{utterance}  
    & & & \multicolumn{4}{p{0.6\linewidth}}{
        \cellcolor[rgb]{0.9,0.9,0.9}{
            \makecell[{{p{\linewidth}}}]{
                \texttt{\tiny{[GM$\langle$P2]}}
                \texttt{ODPOVEĎ: Áno, v treťom stĺpci je poradie presne X, ▢, X, ▢, ▢.} \\
            }
        }
    }
     \\ \\

    \theutterance \stepcounter{utterance}  
    & & & \multicolumn{2}{p{0.3\linewidth}}{
        \cellcolor[rgb]{0.9,0.9,0.9}{
            \makecell[{{p{\linewidth}}}]{
                \texttt{\tiny{[GM$|$GM]}}
                \texttt{continue} \\
            }
        }
    }
    & & \\ \\

    \theutterance \stepcounter{utterance}  
    & & \multicolumn{4}{p{0.6\linewidth}}{
        \cellcolor[rgb]{0.9,0.9,0.9}{
            \makecell[{{p{\linewidth}}}]{
                \texttt{\tiny{[GM$\rangle$P2]}}
                \texttt{Teraz sa rozhodnite. Čo si myslíte: je vaša mriežka a druhá opísaná mriežka rovnaká? Napíšte "ROZHODNUTIE: rovnaká mriežka" ak si myslíte, že ide o rovnakú mriežku, alebo "ROZHODNUTIE: rôzne mriežky" ak si myslíte, že ide o rôzne mriežky. Nepridávajte nič iné.} \\
            }
        }
    }
    & \\ \\

    \theutterance \stepcounter{utterance}  
    & & & \multicolumn{4}{p{0.6\linewidth}}{
        \cellcolor[rgb]{0.9,0.9,0.9}{
            \makecell[{{p{\linewidth}}}]{
                \texttt{\tiny{[GM$\langle$P2]}}
                \texttt{ROZHODNUTIE: rôzne mriežky} \\
            }
        }
    }
     \\ \\

    \theutterance \stepcounter{utterance}  
    & & & \multicolumn{2}{p{0.3\linewidth}}{
        \cellcolor[rgb]{0.9,0.9,0.9}{
            \makecell[{{p{\linewidth}}}]{
                \texttt{\tiny{[GM$|$GM]}}
                \texttt{continue} \\
            }
        }
    }
    & & \\ \\

    \theutterance \stepcounter{utterance}  
    & & & \multicolumn{2}{p{0.3\linewidth}}{
        \cellcolor[rgb]{0.9,0.9,0.9}{
            \makecell[{{p{\linewidth}}}]{
                \texttt{\tiny{[GM$|$GM]}}
                \texttt{success} \\
            }
        }
    }
    & & \\ \\

    \theutterance \stepcounter{utterance}  
    & & \multicolumn{4}{p{0.6\linewidth}}{
        \cellcolor[rgb]{0.9,0.9,0.9}{
            \makecell[{{p{\linewidth}}}]{
                \texttt{\tiny{[P1$\langle$GM]}}
                \texttt{ODPOVEĎ: Áno, v treťom stĺpci je poradie presne X, ▢, X, ▢, ▢.} \\
\texttt{Teraz sa rozhodnite. Čo si myslíte: je vaša mriežka a druhá opísaná mriežka rovnaká? Napíšte "ROZHODNUTIE: rovnaká mriežka" ak si myslíte, že ide o rovnakú mriežku, alebo "ROZHODNUTIE: rôzne mriežky" ak si myslíte, že ide o rôzne mriežky. Nepridávajte nič iné.} \\
            }
        }
    }
    & \\ \\

    \theutterance \stepcounter{utterance}  
    & \multicolumn{4}{p{0.6\linewidth}}{
        \cellcolor[rgb]{0.9,0.9,0.9}{
            \makecell[{{p{\linewidth}}}]{
                \texttt{\tiny{[P1$\rangle$GM]}}
                \texttt{ROZHODNUTIE: rôzne mriežky} \\
            }
        }
    }
    & & \\ \\

    \theutterance \stepcounter{utterance}  
    & & & \multicolumn{2}{p{0.3\linewidth}}{
        \cellcolor[rgb]{0.9,0.9,0.9}{
            \makecell[{{p{\linewidth}}}]{
                \texttt{\tiny{[GM$|$GM]}}
                \texttt{continue} \\
            }
        }
    }
    & & \\ \\

    \theutterance \stepcounter{utterance}  
    & & & \multicolumn{2}{p{0.3\linewidth}}{
        \cellcolor[rgb]{0.9,0.9,0.9}{
            \makecell[{{p{\linewidth}}}]{
                \texttt{\tiny{[GM$|$GM]}}
                \texttt{success} \\
            }
        }
    }
    & & \\ \\

\end{supertabular}
}

\end{document}
